% Suggested replacement for the Scope and limitations paragraph in discussion.tex.
% Values are from abstain_truth_corrected_summary.csv; undefined conditional
% truth rows (zero retained non-malignant truth mass) are excluded explicitly.
\paragraph{Scope and limitations.}
The locked truth is a cell-count share: every typed cell contributes one unit,
irrespective of cell size, RNA content, or capture efficiency.  This is the
estimand supported by the cell-resolved construction, rather than a transcript
mass estimate.  The saved truth objects do not retain a cell-to-spot transcript
total, so an RNA-weighted sensitivity cannot be reconstructed without introducing
an unobserved quantity.  We therefore state the estimand explicitly and avoid
interpreting the reported fractions as RNA abundance.

ABSTAIN changes the estimand for the retained coordinates.  After withholding
the malignant coordinate, the remaining prediction is a conditional composition
and sums to one (maximum absolute closure error $1.1\times10^{-15}$).  In the
corrected truth audit, spots with zero retained non-malignant truth mass were
excluded from conditional error calculations rather than assigned an arbitrary
denominator: 66/400 openST spots, 3/400 realgt spots, and 109/400 realgt3 spots
were undefined under this estimand.  On the defined rows, conditional-neighbor
RMSE was 0.0853 (95\% bootstrap CI 0.0598--0.1111), 0.1682 (0.1409--0.1953),
and 0.1439 (0.1099--0.1785), respectively; corresponding biases were 0.0162,
0.0214, and -0.0004.  These values show why a retained neighbor coordinate
should be read as a conditional share, not as an unbiased estimate of its
original absolute occupancy.  As a descriptive invariant check, the merged
malignant-plus-neighbor mass had RMSE 0.0707, 0.1655, and 0.1545 on the same
spots, because this quantity is unaffected by the conditional renormalization.

The complete scan also shows why pair specification matters.  Historical calls
use biologically designated malignant--neighbor pairs, whereas the scan compares
all represented eligible non-malignant columns and takes their maximum only as a
retrospective identifiability sensitivity.  Under that rule, two of four library
decisions (NSCLC and HCC) differ from the historical designated-pair call; we
therefore retain both estimands and do not present the retrospective maximum as a
pre-registered operating rule.  Neighbor identity remains partly confounded with
library and platform, and the number of complete rows differs between libraries.
The evidence supports an explicit reportability contract and a useful audit of
reference geometry, but it is not a platform-wide calibration study.  Finally,
raw provider-controlled inputs and derived reference materials remain subject to
their data-use terms; the public analysis materials identify locked outputs and
rendering procedures without making a blanket redistribution claim for licensed
inputs.  Clinical outcome inference is outside the scope of this Methods study.
